% Body of the Week 8 study guide (study-guide.tex): the key formulas with
% short explanations, and every exercise of the slides with its detailed
% solution, in the order of the slides. The exercises are in exercises/ (one
% file each); \studyexercise prints the heading, the statement and the
% solution.

\maketitle

\tableofcontents
\clearpage

%======================================================================
\section{Part 1: Permutations and combinations}
%======================================================================

\subsection*{Key formulas}
\addcontentsline{toc}{subsection}{Key formulas}

\begin{key}{Counting toolbox}
	\textbf{Basic principle of counting.} If experiment~1 has $m$ outcomes and,
	for each of them, experiment~2 has $n$ outcomes, then together they have
	$m\cdot n$ outcomes. For $r$ experiments: $n_1n_2\cdots n_r$. Think of a
	tree: the first step gives $m$ branches, each of them splits into $n$
	branches, so the tree has $m\cdot n$ leaves.

	\textbf{Permutations (order matters).} $n$ distinct objects in a row: the
	first place can be filled in $n$ ways, the second in $n-1$ ways, and so
	on, so there are $n! = n(n-1)\cdots 2\cdot1$ orders (and $0!=1$). An
	ordered selection of $k$ out of $n$ (no repetition) stops after $k$
	factors:
	\[ \frac{n!}{(n-k)!}=n(n-1)\cdots(n-k+1). \]
	If repetition is allowed, every place has $n$ choices: $n^k$.

	\textbf{Combinations (order does not matter).} Selection of $k$ out of $n$
	as a set (no repetition):
	\[ \binom nk=\frac{n!}{k!\,(n-k)!}=\frac{n!/(n-k)!}{k!}: \]
	each set comes in $k!$ orders, so it was counted $k!$ times among the
	ordered selections.

	\textbf{Two questions to ask every time.} Does the order matter? Can an
	element be used more than once?
	\begin{center}
		\begin{tabular}{lcc}
			\toprule
			                   & order matters & order does not matter \\\midrule
			no repetition      & $n!/(n-k)!$   & $\binom nk$           \\
			repetition allowed & $n^k$         & $\binom{n+k-1}{k}$    \\\bottomrule
		\end{tabular}
	\end{center}
	The last entry is explained in Part~2 (stars and bars).

	\textbf{Permutations with repetition.} $n$ objects in a row, $n_i$ of them
	identical of kind $i$ ($n_1+\dots+n_r=n$), give
	\[ \frac{n!}{n_1!\,n_2!\cdots n_r!} \]
	different rows: if the identical objects had labels there would be $n!$
	rows, and every row was counted $n_1!\,n_2!\cdots n_r!$ times, once for
	each way of swapping identical objects among themselves
	(Exercise~\exref{letter-arrangements}).

	\textbf{Tricks.} ``Must be together'' $\to$ glue into a block, permute the
	blocks, then permute inside each block. ``Not together'' or ``at least
	one'' $\to$ complement: $\#\text{good}=\#\text{all}-\#\text{bad}$.
\end{key}

%----------------------------------------------------------------------
\studyexercise{license-plates}
\studyexercise{couples-at-a-concert}
\studyexercise{bus-queue}
\studyexercise{letter-arrangements}
\studyexercise{eight-digit-numbers}

%----------------------------------------------------------------------
\subsection*{Multinomial coefficients}
\addcontentsline{toc}{subsection}{Multinomial coefficients}

\begin{key}{Multinomial coefficient}
	The number of ways to divide $n$ distinct objects into $r$ labeled groups
	of sizes $n_1,\dots,n_r$ ($n_1+\dots+n_r=n$) is
	\[ \binom{n}{n_1,\dots,n_r}=\frac{n!}{n_1!\cdots n_r!}
		=\binom{n}{n_1}\binom{n-n_1}{n_2}\cdots\binom{n_r}{n_r}. \]
	\textbf{Why.} Choose the first group in $\binom n{n_1}$ ways, the second
	from the remaining objects in $\binom{n-n_1}{n_2}$ ways, and so on; in the
	product the factorials cancel one after the other.

	\textbf{The same formula as permutations with repetition.} Write under
	every object the name of its group. A division is then a word with $n_1$
	letters ``1'', $n_2$ letters ``2'', \dots, and the number of such words is
	$n!/(n_1!\cdots n_r!)$. With two groups this is the binomial coefficient:
	\[ \binom{n}{k,\,n-k}=\binom nk. \]
	\textbf{Careful.} The formula needs \emph{given} group sizes. If every
	object simply chooses one of $r$ groups, with no condition on the sizes,
	the count is $r^n$ (Exercise~\exref{dividing-into-groups}).
\end{key}

\studyexercise{dividing-into-groups}
\studyexercise{rating-bonds}
\studyexercise{up-and-down-days}
\studyexercise{committees}
\studyexercise{lattice-paths}

\clearpage
%======================================================================
\section{Part 2: Pascal's triangle, sample space, events}
%======================================================================

\subsection*{Pascal's triangle}
\addcontentsline{toc}{subsection}{Pascal's triangle}

\begin{key}{Pascal's triangle}
	\begin{center}
		\begin{tikzpicture}[every node/.style={font=\small}]
			\foreach \n in {0,...,6} {
					\node[navy!70] at (-4.4,-.62*\n) {\scriptsize row $n=\n$};
					\foreach \k in {0,...,\n} {
							\pgfmathtruncatemacro{\c}{factorial(\n)/(factorial(\k)*factorial(\n-\k))}
							\node[inner sep=1pt,minimum size=14pt] (p\n-\k)
							at ({(\k-\n/2)*.95},{-.62*\n}) {\c};
						}
				}
			\draw[gold,thick] (p4-1) circle (7.5pt);
			\draw[gold,thick] (p4-2) circle (7.5pt);
			\draw[teal,thick,fill=teal!15] (p5-2) circle (7.5pt);
			\node[teal] at (p5-2) {10};
			\draw[gold,thick,->] (p4-1)--(p5-2);
			\draw[gold,thick,->] (p4-2)--(p5-2);
			\node[right,align=left,font=\footnotesize] at (3.6,-2.8)
			{$\binom41+\binom42=\binom52$\\ $4+6=10$};
		\end{tikzpicture}
	\end{center}
	Row $n$, position $k=0,1,\dots,n$ (counted from the left, starting at 0)
	holds $\binom nk$. \textbf{Pascal's rule}, proved in
	Exercise~\exref{pascals-rule} below,
	\[ \binom nk+\binom n{k+1}=\binom{n+1}{k+1}, \]
	builds the triangle: every entry is the sum of the two entries just above
	it, e.g.\
	\[ \binom41+\binom42=4+6=10=\binom52. \]
	\textbf{Readings of row $n$:}
	(1) paths with $k$ up-moves among $n$ moves: the lattice paths of
	Exercise~\exref{lattice-paths} and the binomial tree of
	Exercise~\exref{binomial-tree};
	(2) the coefficients of $(a+b)^n$, e.g.\
	\[ (a+b)^5=a^5+5a^4b+10a^3b^2+10a^2b^3+5ab^4+b^5 \]
	(row 5);
	(3) the row sum is $2^n$, and each row is symmetric:
	$\binom nk=\binom n{n-k}$.
\end{key}

\studyexercise{pascals-rule}
\studyexercise{binomial-tree}

%----------------------------------------------------------------------
\subsection*{Binomial and multinomial theorems}
\addcontentsline{toc}{subsection}{Binomial and multinomial theorems}

\begin{key}{Binomial and multinomial theorems}
	\[ (a+b)^n=\sum_{k=0}^n\binom nk a^kb^{n-k},\qquad
		(x_1+\dots+x_r)^n=\sum_{n_1+\dots+n_r=n}\binom{n}{n_1,\dots,n_r}
		x_1^{n_1}\cdots x_r^{n_r}. \]
	\textbf{Why.} Multiply out $(a+b)(a+b)\cdots(a+b)$: from every factor we
	take either $a$ or $b$. We get $a^kb^{n-k}$ when we take $a$ from exactly
	$k$ of the $n$ factors, and there are $\binom nk$ ways to choose these
	factors. This is why the numbers $\binom nk$ are called \emph{binomial
		coefficients}, and why the coefficients of $(a+b)^n$ are row $n$ of
	Pascal's triangle.

	With $r$ terms it is the same: we get $x_1^{n_1}\cdots x_r^{n_r}$ when we
	take $x_1$ from $n_1$ factors, $x_2$ from $n_2$ factors, \dots: a division
	of the $n$ factors into groups of sizes $n_1,\dots,n_r$, counted by the
	multinomial coefficient.

	\textbf{Substitution trick.} The sum of all coefficients of an expansion is
	its value when every variable equals 1. For example
	\[ \sum_{k=0}^n\binom nk=(1+1)^n=2^n,\qquad
		\sum_{k=0}^n(-1)^k\binom nk=(1-1)^n=0\quad(n\ge1). \]
\end{key}

\studyexercise{expansions}

%----------------------------------------------------------------------
\subsection*{Stars and bars}
\addcontentsline{toc}{subsection}{Stars and bars}

\begin{key}{Combinations with repetition (stars and bars)}
	\textbf{Question.} In how many ways can $n$ identical objects be put into
	$r$ distinct boxes?

	\textbf{Example.} 5 identical coins for 3 children. Write a star for each
	coin and a bar between two children:
	\begin{center}
		\begin{tabular}{cc@{\hspace{12mm}}cc}
			\texttt{**|*|**} & means $(2,1,2)$ & \texttt{*****||} & means $(5,0,0)$ \\
			\texttt{|***|**} & means $(0,3,2)$ & \texttt{*||****} & means $(1,0,4)$
		\end{tabular}
	\end{center}
	Every division is a row of 5 stars and 2 bars, and every such row is a
	division. The row has 7 places, and it is fixed once we choose the places
	of the 2 bars: $\binom72=21$ divisions.

	\textbf{In general.} $n$ identical objects into $r$ distinct boxes: $n$
	stars and $r-1$ bars,
	\[ \binom{n+r-1}{r-1}\ \text{ways.} \]
	If every box must be nonempty, first put one object into each box, then
	divide the other $n-r$ freely:
	\[ \binom{(n-r)+r-1}{r-1}=\binom{n-1}{r-1}\ \text{ways.} \]
	\textbf{The last entry of the table.} Choosing $k$ times from $n$ kinds,
	with repetition and without order, is the same problem: only the number
	of times each kind is chosen matters. These are $k$ stars in $n$ boxes,
	so $n-1$ bars:
	\[ \binom{n+k-1}{k} \]
	(Exercise~\exref{identical-dice}).
\end{key}

\studyexercise{stars-and-bars}
\studyexercise{identical-dice}
\begin{minipage}{\textwidth}
	\studyexercise{allocating-capital}
\end{minipage}

%----------------------------------------------------------------------
\subsection*{Sample space and events}
\addcontentsline{toc}{subsection}{Sample space and events}

\begin{key}{Sample space, events and De Morgan's laws}
	\textbf{Sample space} $\Om$: the set of all possible outcomes of an
	experiment. An \textbf{event} is a subset $E\subseteq\Om$; it
	\emph{occurs} if the outcome lies in $E$.
	$E\cup F$ = ``$E$ or $F$'', \ $E\cap F=EF$ = ``$E$ and $F$'', \ $E^c$ =
	``not $E$'', \ $E\subseteq F$ = ``$E$ implies $F$''. $E$ and $F$ are
	\textbf{mutually exclusive} if $EF=\emptyset$.
	\begin{center}
		\setlength{\tabcolsep}{6mm}
		\begin{tabular}{ccc}
			\venn{\fill[teal!35] (0,0) circle (.85);}
			          & \venn{\fill[teal!35] (1,0) circle (.85);}
			          & \venn{\fill[teal!35] (0,0) circle (.85); \fill[white] (1,0) circle (.85);}                       \\
			$E$       & $F$                                                                        & $E\setminus F=EF^c$ \\[2mm]
			\venn{\fill[teal!35] (0,0) circle (.85); \fill[teal!35] (1,0) circle (.85);}
			          & \venn{\begin{scope}\clip (0,0) circle (.85);
					            \fill[teal!35] (1,0) circle (.85);\end{scope}}
			          & \venn{\shadeall{teal!35}\fill[white] (0,0) circle (.85);}                                        \\
			$E\cup F$ & $EF$                                                                       & $E^c$
		\end{tabular}
	\end{center}

	\textbf{De Morgan's laws for two events.}
	\[ \boxed{(E\cup F)^c=E^c\cap F^c}\qquad\qquad
		\boxed{(E\cap F)^c=E^c\cup F^c} \]
	In words: ``not ($E$ or $F$)'' = ``neither $E$ nor $F$''; ``not (both $E$
	and $F$)'' = ``at least one of them fails''.
	\begin{center}
		\renewcommand{\arraystretch}{1.2}
		\begin{tabular}{ccccc}
			\venn{\shadeall{teal!35}\fill[white] (0,0) circle (.85);}
			                  & $\cap$ &
			\venn{\shadeall{gold!40}\fill[white] (1,0) circle (.85);}
			                  & $=$    &
			\venn{\shadeall{navy!35}\fill[white] (0,0) circle (.85);
				\fill[white] (1,0) circle (.85);}       \\
			\scriptsize $E^c$ &        & \scriptsize $F^c$ &  &
			\scriptsize outside both $=(E\cup F)^c$                       \\[1mm]
			\venn{\shadeall{teal!35}\fill[white] (0,0) circle (.85);}
			                  & $\cup$ &
			\venn{\shadeall{gold!40}\fill[white] (1,0) circle (.85);}
			                  & $=$    &
			\venn{\shadeall{navy!35}\begin{scope}\clip (0,0) circle (.85);
					\fill[white] (1,0) circle (.85);\end{scope}} \\
			\scriptsize $E^c$ &        & \scriptsize $F^c$ &  &
			\scriptsize all but the lens $=(E\cap F)^c$
		\end{tabular}
	\end{center}
	Shade $E^c$ and $F^c$: shaded \emph{twice} $=$ intersection; shaded
	\emph{at least once} $=$ union.

	\textbf{General form ($n$ events).}
	\[ \Big(\bigcup_{i=1}^nE_i\Big)^c=\bigcap_{i=1}^nE_i^c,\qquad
		\Big(\bigcap_{i=1}^nE_i\Big)^c=\bigcup_{i=1}^nE_i^c. \]
	Next week we use exactly this: $P(A^cB^c)=1-P(A\cup B)$.

	\textbf{Inclusion--exclusion.} For two events
	\[ |E\cup F|=|E|+|F|-|EF|: \]
	in $|E|+|F|$ the outcomes of the overlap $EF$ were counted twice, so we
	subtract them once. For three events
	\[ |E\cup F\cup G|=|E|+|F|+|G|-|EF|-|EG|-|FG|+|EFG|: \]
	adding the three sets counts the pairwise overlaps twice, so we subtract
	them; but then the center $EFG$ was added three times and subtracted three
	times, so we add it back once. A Venn diagram with three sets is filled
	\emph{from the inside out}: first the center, then the regions of exactly
	two sets, then the regions of exactly one set
	(Exercise~\exref{bank-products}).
\end{key}

\studyexercise{sample-space}
\studyexercise{stock-path-events}
\studyexercise{credit-basket}
\studyexercise{bank-products}
